\chapter{Project Management, Tools, and Impact}
\label{ch:management}

\section{Team Organization}

The project was organized around the four modules described in
Chapter~\ref{ch:design}, with primary ownership assigned per module while
the underlying annotated corpus (Section~\ref{sec:annotation}) was a shared,
jointly-maintained resource that every module depended on.

\begin{table}[H]
\centering
\caption{Module ownership.}
\label{tab:roles}
\begin{tabular}{@{}p{3.4cm}p{9.4cm}@{}}
\toprule
\textbf{Module} & \textbf{Scope} \\
\midrule
A -- Donation-intent classification & Dual-stream cross-attention architecture (v1--v14), binary intent + modifier classification. \\
B -- Persuasion-strategy classification & Taxonomy design, corpus annotation coordination, the 17-iteration multi-label classifier (v1--v17), the decision-rule and ensembling study. \\
C -- Donation-outcome analysis & Logistic outcome modeling, significance testing, single-label projection comparison, temporal-confound analysis. \\
D -- Generic domain extension & Synthetic buyer--seller corpus construction, cross-domain retraining and evaluation of Module A's architecture. \\
\bottomrule
\end{tabular}
\end{table}

\section{Development Methodology}

Module~B in particular followed the iterate-under-a-hard-compute-budget
protocol detailed in Section~\ref{sec:design-methodology}: every one of its
seventeen iterations was smoke-tested locally (GPU-free, against the real
corpus, in under a minute) and then smoke-tested again on Kaggle (a small
data subsample, a few minutes of GPU time) before a full multi-hour training
run was committed, and every full run's results were checked against
deterministic sentinel members before being trusted. This caught real
defects before they cost GPU-hours on several occasions -- most notably when
the category-first reframing (Section~\ref{sec:category-first}) was
implemented, where the local smoke test caught an architectural incompatibility
(a hierarchical-gating step that no longer had a well-defined coarser label
to gate against) before any Kaggle run was attempted.

\section{Tools and Infrastructure}

\begin{table}[H]
\centering
\caption{Core tools and infrastructure.}
\label{tab:tools}
\begin{tabular}{@{}p{3.6cm}p{9.2cm}@{}}
\toprule
\textbf{Purpose} & \textbf{Tool} \\
\midrule
Model training & PyTorch, HuggingFace \texttt{transformers} \\
Shallow baselines, calibration & scikit-learn \\
Statistical / outcome modeling & statsmodels \\
GPU compute & Kaggle (NVIDIA T4$\times$2, free tier, two accounts) \\
Version control & Git \\
Notebook generation & A single builder script regenerates each iteration's notebook, so every change is a reviewable source-code diff rather than a manually edited \texttt{.ipynb} \\
\bottomrule
\end{tabular}
\end{table}

Kaggle's free tier enforces a weekly per-account GPU-hour quota, which a
seventeen-iteration research process is not guaranteed to fit inside a single
account. Every full training run was accordingly queued redundantly across
two independent Kaggle accounts, which both protected the project against
losing a run to quota exhaustion mid-iteration and, on several iterations,
provided an independent reproducibility check by re-running the identical
configuration twice.

\section{Risk Management}

\begin{table}[H]
\centering
\caption{Key risks and mitigations.}
\label{tab:risks}
\begin{tabular}{@{}p{4.3cm}p{8.5cm}@{}}
\toprule
\textbf{Risk} & \textbf{Mitigation} \\
\midrule
GPU quota exhaustion mid-experiment & Redundant training across two Kaggle accounts (Section~\ref{sec:design-methodology}). \\
A multi-hour GPU run fails from a code defect & Mandatory local (GPU-free) and Kaggle-side (data-subsample) smoke tests before every full run. \\
A design change looks like an improvement but is measurement noise & Deterministic sentinel ensemble members required to reproduce exactly before any comparison is trusted (Section~\ref{sec:design-methodology}); decision rules and pruning decisions fitted only on held-out validation data, never the test split. \\
Extreme class imbalance makes the rarest labels unlearnable & Measured, rather than assumed: two strategies with zero test-split support are excluded from macro-F1 by definition (Section~\ref{ch:evaluation}) rather than silently scored as failures; the category-first reframing directly targets the imbalance that \emph{is} addressable. \\
Overclaiming causal or comparative results & Explicit statistical significance testing on every headline comparison (Table~\ref{tab:significance}); the single-label projection control (Section~\ref{sec:phase4-results}) built specifically to avoid an uncontrolled cross-paper comparison; the temporal-confound limitation stated explicitly rather than left implicit (Section~\ref{sec:phase4-results}). \\
\bottomrule
\end{tabular}
\end{table}

\section{Ethical Considerations}

\textbf{LLM-assisted annotation.} A substantial portion of the strategy
annotation was produced by an LLM agent under a human verification
checkpoint (Section~\ref{sec:annotation}), not purely by human annotators.
This is disclosed in the declaration, in the annotation methodology section,
and here, and its reliability is measured directly (Section~\ref{sec:ceiling})
rather than assumed to match human-only annotation quality; every classifier
result in this report is explicitly evaluated against that measured ceiling.

\textbf{Dual-use of persuasion research.} A system that recognizes and
quantifies the effectiveness of persuasion strategies is, by construction,
also a system that could be used to make persuasion more manipulative rather
than more informative -- e.g. to identify which pressure-based tactics
(Table~\ref{tab:odds-ratios} shows \emph{Threat and Pressure} strategies are
associated with a \emph{lower}, not higher, donation rate, which is itself a
finding against using such tactics) extract compliance rather than genuine
willingness. This project's scope is limited to a single benign,
well-studied domain (charitable solicitation for a specific, named charity)
and the associational findings in Chapter~\ref{ch:evaluation} are reported
with their causal limitations stated explicitly (Section~\ref{sec:phase4-results})
specifically so they are not misapplied as a validated manipulation
playbook.

\textbf{Donor and participant privacy.} The underlying P4G corpus consists
of crowd-worker dialogues released publicly by its original authors
\citep{wang2019persuasion} with identifying information already removed;
this project introduces no new personally identifying data and the generic
extension corpus (Module~D) is entirely synthetic.

\textbf{Fairness.} The classifier is deliberately rater-agnostic
(NFR3, Chapter~\ref{ch:requirements}): it is never given the identity of the
turn's annotator as a predictive feature, even though doing so was measured
to increase apparent in-corpus accuracy, because that would be learning a
shortcut specific to this annotation process rather than the underlying
persuasive-language task, and would not transfer to a deployment setting
where the "annotator" of a new turn does not exist.

\section{Broader Impact}

The strategy taxonomy, the annotated corpus, and the strategy-to-outcome
findings (Chapter~\ref{ch:evaluation}) have direct potential use as
fundraiser-coaching material: the odds-ratio analysis
(Table~\ref{tab:odds-ratios}) is, in effect, evidence-based guidance that
reciprocity-and-exchange and information-provision framing are associated
with better outcomes than threat-and-pressure tactics in this domain, for an
organization that wants that guidance grounded in dialogue data rather than
intuition. More broadly, the annotated corpus and the released classifier
give the wider research community a resource for testing persuasion-theory
predictions (Cialdini's principles among them, Chapter~\ref{ch:research})
against real donation-solicitation language at a scale the field has not
previously had for this domain. The project's explicit reporting discipline
-- ceilings, confidence intervals, significance tests, and a stated causal
limitation rather than a suppressed one -- is offered as a template for how
this kind of applied, corpus-driven persuasion research should be reported,
independent of the specific numbers obtained here.
